\documentclass[10pt,a4paper]{amsart}
\usepackage[margin=19mm]{geometry}
\usepackage{amsmath,amssymb,mathtools}
\usepackage[scr=rsfs,cal=euler]{mathalfa}
\usepackage{xfrac}
\usepackage[unicode,hidelinks,bookmarks=false]{hyperref}
\newcommand{\ie}{i.e.\ }
\newcommand{\cf}{cf.\ }
\newcommand{\resp}{respectively\ }
\newcommand{\Subsection}{\subsection}
\providecommand{\nobreakdash}{\nobreak\hbox{-}}
\setlength{\emergencystretch}{2em}
\multlinegap0pt
\usepackage{mathtools}
\usepackage[shortlabels]{enumitem}
\usepackage{amssymb}
\usepackage{xcolor}
\usepackage{tikz}
\newtheorem{theorem}{Theorem}
\newtheorem{lemma}{Lemma}
\newcommand{\la}{\left\langle}
\newcommand{\ra}{\right\rangle}
\newcommand{\trrule}{\mathtt{r}}
\title{On the space of subgroups of Baumslag-Solitar groups II: High transitivity}
\author{Damien Gaboriau, Fran\c{c}ois Le Ma\^itre, Yves Stalder}
\date{Source: arXiv:2410.23224v3 (CC BY 4.0)}
\begin{document}
\maketitle

\noindent\emph{Source and licence.} On the space of subgroups of Baumslag-Solitar groups II: High transitivity. Version arXiv:2410.23224v3 (CC BY 4.0), distributed by arXiv under the Creative Commons Attribution 4.0 International licence, http://creativecommons.org/licenses/by/4.0/, as recorded in the arXiv licence metadata for this exact version and shown on the abstract page https://arxiv.org/abs/2410.23224v3. Full licence notice: \texttt{licenses/arxiv-2410.23224-CC-BY-4.0.txt}.\par\smallskip

\noindent\textit{Editorial scope.} Counted unit: the article's theorem ruled (source0.tex lines 1293--1299). The article's complete original proof of it is delivered with it (source0.tex lines 1390--1436, one \texttt{proof} environment of 47 lines, which the article places away from the statement). No mathematics is written, repaired or re-derived: the statement and the proof are the article's own lines, in the article's own notation, and no retained line is substituted or re-flowed.  Retained uncounted as the source-local context the proof needs: lines 1347--1355.  Nothing was omitted from any retained span, so the package carries no omission record.  No retained line contains a citation, so the wrapper carries no bibliography and no bibliography entry.  No retained line contains a figure environment, an image-inclusion command or drawing code, so no image is delivered and the package declares no asset dependency.  Editorial formatting only, in addition: an AMS \texttt{amsart} preamble that restates the article's own declarations and macros used by the retained lines, namely the environments \texttt{theorem}, \texttt{lemma} on separate counters, together with the standard packages those lines need.  The retained environments print in this document as Theorem 1, Lemma 1 in order of appearance, whereas the article prints the same items with the numbering of its own full document.  The complete native source of the article as distributed in the arXiv e-print for this exact version is bundled unchanged as \texttt{sources/167691/source0.tex}.
	\begin{lemma}
		\label{lem: isomorphisms and one orbit extensions}
		Let $\varphi: X_1 \to X_2$ be an isomorphism between 
		non-saturated pre-actions $\alpha_1 = (X_1,\beta_1, \tau_1)$ to $\alpha_2 = (X_2,\beta_2, \tau_2)$.
		Assume $\alpha'_1$ is a one orbit free \( \trrule \)-extension of $(\alpha_1,x_1)$ and
		$\alpha'_2$ is a one orbit free \( \trrule \)-extension of $(\alpha_2,\varphi(x_1))$
		both positive or both negative.
		Then, $\varphi$ extends to an isomorphism from $\alpha'_1$ to $\alpha'_2$.
	\end{lemma}

	\begin{theorem}\label{thm: forest saturation ruled}
		Let \( \alpha_0 \) be a pre-action and
		let $\trrule$ be a transfer rule. Then 
		\( \alpha_0 \) admits a forest saturation \( \alpha \) satisfying $\trrule$,
		and any two such saturations are isomorphic.
		In particular, up to isomorphism, \( \alpha_0 \) admits a unique maximal forest saturation and a unique minimal forest saturation.
	\end{theorem}

	\begin{proof}[Proof of Theorem \ref{thm: forest saturation ruled}]
		Let \(\alpha_0\) be a non-saturated pre-action on a set \(X\), we first need to show that 
		\(\alpha_0\) admits a saturated forest extension satisfying the transfer rule \(\trrule \). 
		
		First observe that any forest extension of \(\alpha_0\) must have an underlying set
		which is countable or finite if \(X\) was finite, and of the same cardinality as \(X\)
		if \(X\) was infinite. 
		Let us then fix a set \(Z\) containing \(X\) whose cardinality is strictly greater
		than that of \(X\). 
		
		Zorn's lemma provides us a forest extension \(\alpha=(Y,\beta,\tau)\) of \(\alpha_0\) satisfying \(\trrule\) 
		whose underlying 
		set is a subset of \(Z\) and which is maximal among such extensions. 
		Assume by contradiction \(\alpha\) is not saturated, then by our previous observation 
		\(Z\smallsetminus Y\) is infinite and we can thus construct a positive or negative 
		one point free \( \trrule \)-extension of \(\alpha\) with underlying set contained in \(Z\)
		as explained right after the statement of Theorem \ref{thm: forest saturation ruled}.
		This contradicts the maximality of \(\alpha\), which is thus saturated as desired.

		We now prove uniqueness: assume 
		\( \alpha_1, \alpha_2 \) are forest saturations of \( \alpha_0 \) satisfying \( \trrule \).
		For \( i=0,1,2 \), write \( \alpha_i = (X_i, \beta_i, \tau_i) \).
		Consider the set of partial isomorphisms,
		from a forest extension \( \alpha_1' = (X_1',\beta_1',\tau_1') \) contained in \( \alpha_1 \) 
		to a forest extension \( \alpha_2' = (X_2',\beta_2',\tau_2') \) contained in \( \alpha_2 \),
		that fix \( X_0 \) pointwise.
		This is a non-empty inductive poset.
		Hence, it contains a maximal element \( \varphi:X_1' \to X_2' \) by Zorn's Lemma.
		
		Let us prove by contradiction that $\alpha_1' = \alpha_1$ and $\alpha_2' = \alpha_2$; this will complete the proof.
		If this is  not the case, we may assume that $\alpha_1' \neq \alpha_1$ 
		(the case $\alpha_2' \neq \alpha_2$ is similar).
		In this case, $\alpha_1'$ is non-saturated,
		hence one can find \(x\in X'_1\) and a restriction \(\alpha''_1\) of \( \alpha_1 \)
		which is a one-orbit free \( \trrule \)-extension of \( (\alpha_1',x_1) \).
		As $\varphi: X_1' \to X_2'$ is an isomorphism
		the orbit $\varphi(x) \la \beta_2 \ra$ lies in $X_2'$
		while $\varphi(x) \tau_2^\varepsilon \la \beta_2 \ra$ does not.
		So, the restrictions of $\beta_2, \tau_2$ to 	
		\(X_2'' \coloneqq X_2' \sqcup \varphi(x) \tau_2^\varepsilon \la \beta_2 \ra\)
		give a one-orbit free \(\trrule\)-extension $\alpha_2''$ of $\alpha_2'$, contained in $\alpha_2$.
		Notice that $x \la \beta_1 \ra$ and $\varphi(x) \la \beta_2 \ra$
		have the same cardinal since $\varphi$ is an isomorphism.
		Lemma \ref{lem: isomorphisms and one orbit extensions} now implies
		that $\varphi$ extends to an isomorphism from $\alpha_1''$ to $\alpha_2''$, 
		contradicting the maximality of $\varphi$.
	\end{proof}

\end{document}
